\subsection{All balanced classes vanish in the fixed two-prime tower}

\begin{statement}{A.179}{Proved}
{Cyclic pure class groups and one coherent completed norm family}
Retain the actual compatible roots and field inclusions
\[
 \begin{gathered}
 \rho_n^{3^n}=3\cdot649^{24},\qquad
 \rho_{n+1}^3=\rho_n,\\
 K_n=\Q(\rho_n),\qquad E_n=K_n(\zeta_3),
 \qquad S=\{2,3,11,59,\infty\}.
 \end{gathered}
\tag{A.179.1}
\]
For every \(n\ge2\), the three-primary ordinary class group
\(\operatorname{Cl}(K_n)\{3\}\) is cyclic and is generated by
\[
 R_{59,n}=[P_{59,2,n}]_{(3)}-2[P_{59,0,n}]_{(3)}.
\tag{A.179.2}
\]
The labels are the residue-degree-two and residue-degree-one
changed-prime branches of A.155.  This assertion specifies a
generator, without asserting an order at the previously uncomputed
levels.  In particular,
\[
                         C_n^{\rm bal}=0
                         \qquad(n\ge1).
\tag{A.179.3}
\]
There exists a compatible family in the stated pro-three completions,
\[
 b_n\in\widehat{\mathcal O_{K_n,S_n}^{\times}}_{(3)},\qquad
 N_{K_{n+1}/K_n}b_{n+1}=b_n,\qquad
 N_{K_n/\Q}b_n=[-2]\qquad(n\ge1).
\tag{A.179.4}
\]
Thus the fixed three-primary one-variable obstruction of A.77,
A.81 and A.155 for this radicand and target is zero, with its actual
norm transitions retained.  At every finite height there is also
a chain of actual \(S\)-units with rational norm exactly \(-2\).
The infinite-family assertion in A.179.4 is an assertion about
pro-three completions.  No inverse limit of the discrete actual
unit groups is asserted, and its first two coordinates are not
prescribed to equal the particular representatives in A.176.
\end{statement}

\begin{proof}
\emph{The class maps and the persistent anti-invariant detectors.}
Write
\[
 A_n=\operatorname{Cl}(E_n)\{3\},\qquad
 \tau(\zeta_3)=\zeta_3^{-1},\quad \tau|_{K_n}=1.
\tag{A.179.5}
\]
A.177 proves that \(A_n\) is generated by the six prime classes
above eleven and fifty-nine, for every \(n\ge2\).
Every one of these primes is totally ramified in
\(E_{n+1}/E_n\).  Its upper prime ideal is therefore fixed by
\(\sigma_n\), a generator of this cyclic cubic extension.
It follows on these actual generators that
\[
                    \sigma_n|_{A_{n+1}}=1.
\tag{A.179.6}
\]
Ideal norm \(A_{n+1}\to A_n\) is onto: each of the six generating
lower primes is the norm of its upper prime, with residue-degree
weight one.  Ideal extension \(j_n\) and norm satisfy
\(j_nN_n=1+\sigma_n+\sigma_n^2=3\) on \(A_{n+1}\).
These are the literal maps of the given tower.

For each changed prime use the labels of A.173, and put
\(\Delta_{q,n}=[R_{q,1,n}]_{(3)}-[R_{q,2,n}]_{(3)}\).
The involution exchanges these two primes and fixes the third.
The two unramified cubic Kummer characters of A.173 are given by
\[
                       z_1=11(1+6\alpha),\qquad z_2=649.
\tag{A.179.7}
\]
Their base changes to every \(E_n\) are unramified.  Denote the
resulting ordinary class characters by \(\chi_{1,n},\chi_{2,n}\).
Restriction of an unramified character sends upper arithmetic
Frobenius to lower arithmetic Frobenius raised to the residue degree.
The successive changed-prime residue degrees are one, so the exact
pairing of A.173 persists:
\[
 \begin{pmatrix}
 \chi_{1,n}(\Delta_{11,n})&\chi_{1,n}(\Delta_{59,n})\\
 \chi_{2,n}(\Delta_{11,n})&\chi_{2,n}(\Delta_{59,n})
 \end{pmatrix}
                 =\begin{pmatrix}0&2\\2&1\end{pmatrix}
                 \quad\text{over }\F_3.
\tag{A.179.8}
\]
This also proves that neither required character is lost under
base change.  In particular the two indicated anti-invariant prime
classes are linearly independent in \(A_n/3A_n\) for every \(n\).

\emph{Descent of an arbitrary plus cubic class character.}
Let
\(\lambda:A_{n+1}\to\F_3\) be nonzero and satisfy
\(\lambda\tau=\lambda\).
Class field theory associates an everywhere unramified cyclic
cubic extension \(M/E_{n+1}\) to its kernel, with the indicated
identification of its Galois group with \(\F_3\).
Equation A.179.6 makes this kernel stable under \(\sigma_n\);
conjugation acts trivially on the quotient because \(\lambda\)
itself is invariant.  Consequently \(M/E_n\) is Galois, with
a central exact sequence
\[
 0\longrightarrow \F_3
   \longrightarrow G=\operatorname{Gal}(M/E_n)
   \longrightarrow\operatorname{Gal}(E_{n+1}/E_n)
   \longrightarrow0.
\tag{A.179.9}
\]
Thus \(G\) is one of the two abelian groups of order nine.

There is a unique prime above three in \(E_n\) and in
\(E_{n+1}\), and its latter extension is totally ramified of
degree three.  Since \(M/E_{n+1}\) is unramified, the inertia
subgroup \(I\) of \(G\) above three has order three and maps
isomorphically onto \(\operatorname{Gal}(E_{n+1}/E_n)\).
The cyclic group of order nine has only one order-three subgroup,
namely the kernel in A.179.9, and cannot have this property.
Therefore
\[
                            G\simeq\F_3^2.
\tag{A.179.10}
\]
This argument uses the actual ramification indices, not a choice
of an abstract splitting of A.179.9.

The automorphism \(\tau\) also preserves the kernel of the class
character.  It therefore acts by conjugation on \(G\).
This action is independent of its lift, because \(G\) is abelian,
and has square one.  It is the identity on the kernel in A.179.9,
by \(\lambda\tau=\lambda\), and is multiplication by minus one
on the quotient: on the actual cubic root extension,
\(\tau\sigma_n\tau^{-1}=\sigma_n^{-1}\).
As two is invertible over \(\F_3\), \(G\) has a plus kernel line
and a minus complementary line.  The inertia subgroup \(I\) is
stable under \(\tau\): in an abelian extension all primes above
a fixed lower prime have the same inertia subgroup, and the lower
three-prime is fixed.  Its projection is the minus quotient, so
it is precisely that minus line.

Let \(D=M^I\).  Then \(D/E_n\) is cyclic cubic, \(\tau\) acts
trivially on its Galois group, and
\[
                             D E_{n+1}=M.
\tag{A.179.11}
\]
The Kummer correspondence over \(E_n\), which contains \(\mu_3\),
identifies this extension with a nonzero class
\(d\in E_n^\times/E_n^{\times3}\).
Its equivariance is the perfect pairing
\(\operatorname{Gal}(D/E_n)\times\langle d\rangle\to\mu_3\).
The first factor is \(\tau\)-fixed and the target is inverted;
therefore
\[
                               \tau d=-d
                 \quad\text{in }E_n^\times/E_n^{\times3},
\tag{A.179.12}
\]
where the displayed group is written additively.  Equivalently a
representative satisfies \(\tau(d)d\in E_n^{\times3}\).
This proves the precise twist in the radical descent.

\emph{The descended radical is already everywhere unramified.}
Outside three, eleven and fifty-nine, the extension
\(E_{n+1}/E_n\) is unramified.  Unramified base change does not
change the valuation of a radical.  Since the extension defined
by \(d\) is unramified upstairs, all these lower valuations are
divisible by three.  This includes every prime above two.
At three, the unique prime is fixed by \(\tau\), so A.179.12
gives \(2v_3(d)=0\) modulo three and again
\(v_3(d)=0\) modulo three.

At each changed prime the fixed branch has valuation zero modulo
three, and the two exchanged branches have opposite valuations,
by the same equivariance.  For integers \(k_{11},k_{59}\), the
actual principal divisor of a representative consequently has
the form
\[
 \operatorname{div}(d)=3D_0+
 k_{11}(R_{11,1,n}-R_{11,2,n})+
 k_{59}(R_{59,1,n}-R_{59,2,n}),
\tag{A.179.13}
\]
with an actual finite fractional ideal divisor \(D_0\).
In \(A_n/3A_n\), this is a relation between the two classes in
A.179.8.  Evaluating both characters there gives
\(k_{11}=k_{59}=0\) modulo three.
Every finite valuation of \(d\) is therefore divisible by three.
At residue characteristic different from three, a unit cubic
Kummer extension is unramified, and removing a cube of a
uniformizer reduces to that case.  Thus the descended extension
is already unramified at every place other than three.

For clarity, the local comparison at three retains its actual
radical line.  Put \(k=(E_n)_3\) and
\(l=k(\sqrt[3]{\rho_n})=(E_{n+1})_3\).
The extension is totally ramified cyclic cubic and
\(v_k(\rho_n)=2\).  Let \(k^{\rm ur,3}/k\) be its unramified
cubic extension.  The full inverse image of the unramified
Kummer line under restriction to \(l\) is
\[
 \langle[\rho_n]\rangle+
 \langle\text{the radical of }k^{\rm ur,3}/k\rangle
          \ \subset k^\times/k^{\times3}.
\tag{A.179.14}
\]
Indeed, a cubic extension \(k'/k\) whose compositum with \(l\)
is unramified over \(l\) lies in \(l k^{\rm ur,3}\): a local
field has a unique unramified extension of any given degree.
Conversely every cubic subfield of that elementary abelian
compositum has the required property.  This proves both
directions of the displayed equality, including the split case.
An unramified radical has valuation divisible by three, since
valuation not divisible by three gives a totally ramified cubic.
The valuation of \(\rho_n\) is two, so the already established
valuation condition on \(d\) excludes its component along
\([\rho_n]\).  The extension \(D/E_n\) is unramified at three
as well.  All infinite places of \(E_n\) are complex.

It follows that \(D/E_n\) is an ordinary class field.  Its
\(\tau\)-invariant class character pulls back to \(\lambda\)
by A.179.11, after choosing the generator of its cyclic Galois
group to agree with the given generator upstairs.
We have proved surjectivity of the actual pullback map
\[
 \operatorname{Hom}(A_n^+,\F_3)
       \longrightarrow\operatorname{Hom}(A_{n+1}^+,\F_3).
\tag{A.179.15}
\]
The map is also injective: it is dual to ideal norm, which is
surjective on the plus summands by A.179.6 and the idempotent
\((1+\tau)/2\).  Equivalently an unramified cubic extension
cannot be the totally ramified extension \(E_{n+1}/E_n\).

\emph{The pure class group and the balanced quotient.}
A.173 gives \(A_2^+\simeq\Z/3\).
Induction in A.179.15 shows that every \(A_n^+\) has
three-rank one; being a finite abelian three-group, it is cyclic.
For the actual quadratic extension \(E_n/K_n\), ideal extension
\(j\) and norm satisfy \(Nj=2\) and \(jN=1+\tau\), including
ramified ideal multiplicities.  Since two is invertible on the
three-primary groups, these identities identify
\(\operatorname{Cl}(K_n)\{3\}\) with \(A_n^+\).

The explicit compatible character \(\chi_n\) of A.170 has
values zero on \(P_{59,2,n}\) and one on \(P_{59,0,n}\).
Thus
\[
                             \chi_n(R_{59,n})=1.
\tag{A.179.16}
\]
In a cyclic three-group an element detected by a cubic character
generates the group: its coordinate is not divisible by three.
This proves A.179.2.  The generator is one of the literal
relations divided out in the definition A.155.2 of
\(C_n^{\rm bal}\), so that entire quotient is zero.
At level one the class group is trivial by the class-number-one
proof in A.165.  This proves A.179.3 at every level.

\emph{One completed family and actual finite chains.}
For each \(n\), the unit group \(\mathcal O_{K_n,S_n}^{\times}\)
is finitely generated.  Its pro-three completion is compact.
The exact criterion A.155.7 and A.179.3 say that its fibre of
norm \([-2]\) is nonempty.  Norm maps are continuous and are
the given field norms.  Any finite list of compatibility
equations is solved by choosing a preimage at its largest level
and taking successive norms.  The closed equations consequently
have the finite intersection property in the product of the
compact completed unit groups.  Their common intersection gives
A.179.4.  A.166 identifies the same family inside the unrestricted
completed multiplicative norm fibre; in particular support has
not been replaced by a new hypothesis in this comparison.
The maps from these actual root-torsor norm quotients to the
one-variable completed coefficient module are those of A.77
and A.81, so the compatible family kills that connecting class.

There is also a direct finite-height passage to actual elements.
Set \(U_0=\mathcal O_{\Q,S}^{\times}\) and
\(U_n=\mathcal O_{K_n,S_n}^{\times}\).
The quotient
\[
                           U_0/N_{K_n/\Q}U_n
\tag{A.179.17}
\]
is a finite group killed by \(3^n\): every \(u\in U_0\),
viewed as a scalar in \(K_n\), has norm \(u^{3^n}\).
Completion of finitely generated abelian groups is exact, and
does not change this finite three-group.
The completed preimage of \([-2]\) therefore says exactly that
the actual class of \(-2\) in A.179.17 is zero.
There is an actual \(S_n\)-unit of norm \(-2\) at each finite
height.  Its successive actual norms give the asserted finite
chain.  The compactness argument above applies to the completed
groups and asserts no compactness of their discrete actual
subgroups.

\emph{Finite checks.}
The exact coefficient, inertia and eigenspace calculations are
retained in \path{code/two_inert_all_height_balance_certificate.py},
with its executed output.  They check the central order-nine
descent, the twisted Kummer signs, the two-by-two detecting matrix,
the valuation-line exclusion and the cyclic-generator deduction.
The uniform class-field and compactness arguments above supply
the all-height statements.
\end{proof}
